%专门用于幻灯片制作的beamer文档类
\documentclass[xcolor=table]{beamer}
%在beamer中使用中文
\usepackage[UTF8]{ctex}


%设置主题
\usetheme{CambridgeUS}
\usecolortheme{dolphin}%蓝色

%调用arev宏包修改字体
\usebeamerfont{professionalfonts}
\usepackage{arev}


%列表
\usepackage{enumerate}%有序列表
\usepackage{pifont}%解说列表description


%表格
\usepackage{graphicx}%表格
\usepackage{booktabs, multicol, multirow}%调整表格行高
\usepackage{bigstrut}


%文本框
\usepackage{tcolorbox}
\tcbuselibrary{skins,vignette,raster,listings,listingsutf8,theorems, breakable,magazine,fitting}


%图像处理
\usepackage{graphicx}
\usepackage{float}%允许浮动
\usepackage{caption}%设置标题
\usepackage{graphicx, subfig}


%数学公式
\usepackage{mathtools}%花括号括起的公式


%Tikz绘图
\usepackage{tikz}
\usepackage{pgfplots}



%定理环境
%格式和字体
%\usepackage{theorem}%对定理类环境的进行设置
%\usepackage{ntheorem}%拓展theorem宏包的功能
%\theoremstyle{plain}%定理类环境的格式
%\theorembodyfont{}%定理体字体
%\theoremheaderfont{\heiti}%定理头字体

%自定义定理环境
\newtheorem{defi}{\textbf{定义~}}[section]
\newtheorem{exam}{\textbf{例~}}[section]
\newtheorem{thm}{\textbf{定理~}}[section]
\newtheorem{lem}{\textbf{引理~}}[section]
\newtheorem{cor}{\textbf{推论~}}[section]
\newtheorem{prop}{\textbf{命题~}}[section]
\newtheorem{rem}{\textbf{注~}}[section]


\renewcommand\proofname{\textbf{证明:}}


%自定义宏
\newcommand{\normalfun}[1]
{exp(-x^2/#1)}

\newcommand{\drawfun}[2]
{
		\addplot[
		domain=-3:3,
		samples=200,
		color=#1,
		]
		{#2};
}